\scanpage{20}
Ceci dit, on a des isom. can.
\[
\left\{\begin{aligned}
\unsure{\Phi|\mathscr L|\mathcal L}_A&=({}_pA^*)_{\mathrm{\acute et}}^\vee\\
\Phi_{A^*}&=({}_pA)_{\mathrm{\acute et}}^\vee
\end{aligned}\right.
\]

\uline{NB.} Soit $\varphi$ une polarisation de $A$ (plus
gén. une \deleted{\unsure{diviseur}} section de $\underline{\mathrm{NS}}_{A/S}$). Alors
$\varphi$ définit une \deleted{forme alternée sur}\inserted{section \deleted{\unsure{dite de}} $\varphi'$ de}
$\bigwedge^2\mathrm{DR}^1(A/S)$, \deleted{\unsure{à} valeurs en $\underline O_S$. Cette forme admet
les composantes $\underline H^{10}=\underline t_A^\vee$ et $H^{01}=\underline t_{A^*}$ comme
sous-espaces isotropes, et définit donc un
accouplement $\underline t_A^\vee\otimes\underline t_{A^*}\mathbin{\to}\underline O_S$ i.e.
un hom $\underline t_{A^*}\longrightarrow\underline t_A$. Cet hom n'est
autre que l'hom tangent à $\widetilde\varphi:A\mathbin{-}A^*$}
\[
\bigwedge^{2}\mathrm{DR}^1(A/S)\simeq\bigwedge^2\underline t_A^\vee+
\overbrace{\underline t_A^\vee\otimes\underline t_{A^*}}^{\operatorname{Hom}(\underline t_A,\underline t_{A^*})}
+\bigwedge^{\unsure{2}}\underline t_{A^*}
\]
Ceci dit, $\varphi'$ est dans le terme médian, et
comme tel n'est autre que l'hom
tangent $\widetilde\varphi':\underline t_A\mathbin{-}\underline t_{A^*}$ de $\widetilde\varphi:A\to A^*$ défini
par $\varphi$, qui est déterminé \unsure{également|ipso facto|formellement}, par tensorisation
\deleted{par} $\otimes_{\mathbb F_p}\underline O_S$, par
\diagram{\begin{tikzcd}[ampersand replacement=\&]
\underline\Phi_A \arrow[r,dash] \arrow[d,phantom,"\shortparallel" description] \& \underline\Phi_{A^*} \arrow[d,phantom,"\shortparallel" description]\\
({}_pA^*)_{\mathrm{\acute et}}^\vee \arrow[r,dash] \& ({}_pA)_{\mathrm{\acute et}}^\vee
\end{tikzcd}} transposé de
$({}_pA)_{\mathrm{\acute et}}\mathbin{-}({}_pA^*)_{\mathrm{\acute et}}$ induit par $\,{}_p\widetilde\varphi:{}_pA\mathbin{-}{}_pA^*$ \ldots
% archive sheet number: 19
